\chapter{Results II --- Cross-testbed transfer and within-HAI LOPO}
\label{ch:results-transfer}

This is the headline chapter.  Notebook: \file{05\_cross\_dataset.ipynb}.  Figures: F2, F3, F4 under \file{results/figures/05\_cross\_dataset/}.

\section{Schema alignment}

Six canonical types intersect HAI and Morris after hand-tagging in \file{data/feature\_types.yaml}: pressure, pump\_state, setpoint, system\_state, valve\_position, control\_signal.  HAI contributes 79 features collapsing to 6; Morris 16 features collapse to 6.  The projection is lossy by design --- we trade per-variable fidelity for a cross-testbed input shape that every detector can consume.

\section{Finding 3.1 --- source calibration collapses to the class-prior baseline}
\label{sec:finding-src-cal}

Under \code{src-cal} (threshold $= $ 99th percentile of source validation scores), every transferred detector degenerates to a predict-all-attack baseline calibrated to the \emph{target's} class prior.

\begin{table}[h]
\centering
\caption{Source-calibrated cross-testbed F1 tracks the target class prior rather than any detector property.}
\small
\begin{tabular}{lrrl}
\toprule
Target & Class prior (windowed) & Predict-all F1 & Observed F1 range \\
\midrule
Morris & 0.91 & 0.95 & 0.95--0.99 (windowed); 0.62--0.75 (classical) \\
HAI    & 0.03 & 0.05 & 0.04--0.07 (all models) \\
\bottomrule
\end{tabular}
\end{table}

The score distributions on source and target data have no calibrated overlap; a fixed percentile threshold either flags everything or nothing.  \textbf{A paper reporting only \code{src-cal} numbers on Morris (say F1$=0.98$) is reporting a class-prior artifact, not a transfer success.}

\begin{figure}[h]
\centering
\includegraphics[width=0.95\textwidth]{05_cross_dataset/transfer_same_vs_cross.png}
\caption{Same-domain target F1 vs.\ source-calibrated transfer F1 per model.  Transfer numbers mirror the target's predict-all baseline (grey); the observed F1 is not detecting anything.}
\label{fig:transfer-same-vs-cross}
\end{figure}

\section{Finding 3.2 --- target calibration reveals a classical-vs-SOTA split}
\label{sec:finding-tgt-cal}

Under \code{tgt-cal} (threshold from 99th percentile of \emph{target}-normal validation --- unlabeled, attack-free data available in any realistic deployment), the collapse disappears and a genuine separation emerges.

\begin{table}[h]
\centering
\caption{HAI $\rightarrow$ Morris transfer F1 under \code{tgt-cal}.  Windowed SOTA models separate from classical ones only once thresholds are target-calibrated.}
\small
\begin{tabular}{lrr}
\toprule
Model & tgt-cal eTaPR & tgt-cal PA-F1 \\
\midrule
Dense AE          & 0.07 & 0.14 \\
OC-SVM            & 0.08 & 0.17 \\
Isolation Forest  & 0.13 & 0.22 \\
\textbf{LSTM-AE}  & \textbf{0.44} & \textbf{0.81} \\
\textbf{TranAD}   & \textbf{0.38} & \textbf{0.77} \\
\textbf{USAD}     & \textbf{0.40} & \textbf{0.74} \\
\bottomrule
\end{tabular}
\end{table}

The classical three are clustered near zero; the windowed SOTA three sit at $0.38$--$0.44$ eTaPR.  This is the first crisp classical-vs-SOTA gap in our study and it is invisible under source calibration.

The reverse direction (Morris $\rightarrow$ HAI) is harder: pointwise F1 stays under $0.07$ for all models because the HAI $2.4\%$ attack prior saturates any 99th-percentile operating point.  PA-F1 recovers for classical detectors (Dense AE $0.47$, OC-SVM $0.36$) via the segment-inflation mechanism critiqued in Chapter~\ref{ch:results-detection} --- the same weakness that boosts classical PA-F1 same-dataset.

\begin{figure}[h]
\centering
\includegraphics[width=0.98\textwidth]{05_cross_dataset/transfer_calibration_comparison.png}
\caption{Source-cal vs.\ target-cal transfer F1 vs.\ same-domain baseline for each model.  Target calibration alone converts the trivial-baseline collapse into a meaningful measurement.}
\label{fig:transfer-cal-comp}
\end{figure}

\subsection{Why only the windowed models transfer}

Hypothesis: the classical models' normal manifold is essentially per-dataset specific; they fit the density of the source type-vector without capturing the temporal structure that is approximately shared between testbeds.  Windowed models encode short-term dynamics of normal operation --- pressure rising in ramp phases, valves opening after setpoints move --- and those dynamics recur in both testbeds.  The representation retains the invariant; the operating point does not.

\section{Finding 3.3 --- within-HAI LOPO isolates feature shift}
\label{sec:finding-lopo}

Holding class prior fixed (both train and test are HAI), we drop one process's feature columns and compute F1 per-attack-target-process.

\subsection{P2 drop produces targeted blindness}

\begin{table}[h]
\centering
\caption{F1 on P1-attacks vs.\ P2-attacks when P2 features are held out.  The drop is targeted to P2, not global.}
\small
\begin{tabular}{lrr}
\toprule
Model & F1 on P1-attacks (drop-P2) & F1 on P2-attacks (drop-P2) \\
\midrule
Dense AE & 0.41 & 0.03 \\
LSTM-AE  & 0.09 & 0.00 \\
USAD     & 0.28 & 0.07 \\
\bottomrule
\end{tabular}
\end{table}

P2 is the turbine controller; its sensors are almost all discrete state flags with no analogue in P1, P3, or P4.  Removing them leaves the model blind to attacks that manifest only through those flags --- while detection on P1 attacks barely moves.

\subsection{P1 drop produces global degradation}

When P1 (water loop, $33$ of $79$ features, $82\%$ of test attacks) is removed, all three models drop to F1 $\sim 0.1$ on every attack class.  P1 carries most of the model's picture of ``normal''; losing those features starves the whole detector.

\subsection{P3 and P4 drop is mild and sometimes helpful}

Dropping P3 or P4 occasionally \emph{improves} detection: USAD drop-P3 F1 on P1-attacks $= 0.57$ (vs.\ $0.35$ in the Phase-2 same-domain baseline).  Narrower input regularizes the windowed models.

\begin{figure}[h]
\centering
\includegraphics[width=0.95\textwidth]{05_cross_dataset/lopo_heatmap.png}
\caption{LOPO per-attack-process pointwise F1.  Rows are the held-out process; columns are the process the attack targets.  P4 column is empty because HAI 21.03 has no P4 attacks.}
\label{fig:lopo-heatmap}
\end{figure}

\subsection{No consistent SOTA advantage under LOPO}

Unlike the cross-testbed setting (where SOTA clearly beat classical under \code{tgt-cal}), LOPO shows no consistent classical/SOTA gap.  Dense AE is competitive with USAD on the P2-blindness diagnostic; both beat LSTM-AE on absolute numbers under drop-P2.  The cross-testbed SOTA advantage therefore looks tied to \emph{cross-schema representation}, not to robustness to feature loss per se.

\section{Cross-reading: sensor-level transfer ceilings}

Best within-HAI LOPO pointwise F1 $= 0.57$ (USAD, drop-P3, P1-attacks).  Best cross-testbed \code{tgt-cal} pointwise F1 $\approx 0.13$ (LSTM-AE HAI$\rightarrow$Morris).  Same architecture, same class-prior control, partial sensor overlap: the cross-dataset ceiling is an order of magnitude lower than the within-dataset one.  This defends the cross-testbed numbers against the ``you just need a better common feature space'' reviewer critique --- even under the best within-testbed conditions, sensor-level transfer loses half its detection capability.

\section{External-validation pipeline}
\label{sec:external-validation}

The transfer study above is closed over the two datasets chosen in Chapter~\ref{ch:method}.  To allow any later Morris-family capture (re-collected gas-pipeline ARFFs, the water-storage companion, future variants) to be scored against the same trained detectors without retraining, we ship a small external-validation pipeline: each run in \file{experiments/configs/} can opt into saving its fitted model, scaler, and threshold as a self-contained artifact, and a thin CLI (\code{python -m scripts.score\_external}) plus FastAPI app (\file{app/main.py}) apply that artifact to any uploaded ARFF/CSV.  The artifact reproduces the original summary-parquet row to floating-point tolerance on the held-out test split, so the pipeline is suitable both for the reproducibility gate (verifying no drift between train-time metrics and post-hoc scoring) and for extending the cross-testbed analysis beyond what is in this thesis.  The tooling is strictly Phase~A: gas-pipeline adapter only, localhost only, labels required for the metric panel; water-storage and power-system adapters are future work (\S\ref{sec:future}).
